\documentclass[a4paper]{article}

\usepackage[utf8x]{inputenc}    
\usepackage[T1]{fontenc}
\usepackage[spanish]{babel}
\usepackage{multicol}

\usepackage{wrapfig}
\usepackage{graphicx}

\usepackage{bm}
\usepackage{amsxtra} 
\usepackage{amssymb}% to get the \mathbb alphabet
\usepackage{amsmath}

\usepackage[box,completemulti,separateanswersheet]{automultiplechoice}    
\def\AMCformQuestion#1{\vspace{\AMCformVSpace}\par {\sc Pregunta #1:} }    
\def\AMCbeginQuestion#1#2{\par\noindent{\bf Pregunta #1}#2\hspace*{1em}}
\def\AMCcleardoublepage{\ifodd\thepage\clearpage\mbox{}\fi\clearpage}

\begin{document}

\AMCrandomseed{1237893}

%%%%%%%%%%%%%%%%%%%%%%%%%%%%%%%%%%%%%%%%%%%%%%%%%%%%%%%%%%%%%%%%%%%%%%%%%%%%%
\element{test1}{
\begin{question}{p1}
En la formulación de elementos finitos con integración selectiva-reducida:
\begin{choices}
	\correctchoice{La matriz de rigidez se descompone en una parte
volumétrica, que se subintegra, y otra parte desviadora con integración
completa.}
	\wrongchoice{Las matrices de rigidez de los elementos con cargas aplicadas se subintegran.}
	\wrongchoice{En una parte reducida de la malla, las matrices de
rigidez elementales se integran con una cuadratura selectiva.}
	\wrongchoice{Ninguna de las respuestas anteriores es correcta.}
\end{choices}
\end{question}
}
%%%%%%%%%%%%%%%%%%%%%%%%%%%%%%%%%%%%%%%%%%%%%%%%%%%%%%%%%%%%%%%%%%%%%%%%%%%%%
\element{test1}{
\begin{question}{p2}
La formulación $\overline{\bm{B}}$ de elementos finitos:
\begin{multicols}{2}
\begin{choices}
	\correctchoice{Se basa en modificar la matriz de interpolación de
deformaciones estándar $\bm{B}$}
	\wrongchoice{Se basa en remplazar la matriz de interpolación de
deformaciones por un vector de deformaciones libre de bloqueo}
	\wrongchoice{Evita el problema de bloqueo por flexión pero no el
bloqueo por incompresibilidad}
	\wrongchoice{Evita el problema de bloqueo, pero no pasa la prueba de
la parcela}
\end{choices}
\end{multicols}
\end{question}
}
%%%%%%%%%%%%%%%%%%%%%%%%%%%%%%%%%%%%%%%%%%%%%%%%%%%%%%%%%%%%%%%%%%%%%%%%%%%%%
\element{test1}{
\begin{question}{p3}
Los elementos con formulación de modos incompatibles:
\begin{multicols}{2}
\begin{choices}
	\correctchoice{Incorporan funciones de forma asociadas a grados de libertad internos (los desplazamientos asociados a las mismas son nulos en los
nodos del elemento)}
	\wrongchoice{Al igual que ciertos elementos mixtos, tienen grados
de libertad de presión}
	\wrongchoice{Son una clase de elementos isoparamétricos}
	\wrongchoice{Ninguna de las respuestas anteriores es correcta}
\end{choices}
\end{multicols}
\end{question}
}
%%%%%%%%%%%%%%%%%%%%%%%%%%%%%%%%%%%%%%%%%%%%%%%%%%%%%%%%%%%%%%%%%%%%%%%%%%%%%
\element{test1}{
\begin{question}{p4}
El valor del desplazamiento en dirección $X$ del nodo $A$,
calculado con los elementos isoparamétricos, vale aproximadamente:
\begin{multicols}{2}
\begin{choices}
	\correctchoice{$-1.7$ mm}
	\wrongchoice{$-2.8$ mm}
	\wrongchoice{$-0.7$ mm}
	\wrongchoice{$-0.2$ mm}
\end{choices}
\end{multicols}
\end{question}
}
%%%%%%%%%%%%%%%%%%%%%%%%%%%%%%%%%%%%%%%%%%%%%%%%%%%%%%%%%%%%%%%%%%%%%%%%%%%%%
\element{test1}{
\begin{question}{p5}
El valor del desplazamiento en dirección $X$ del nodo $A$,
calculado con los elementos mixtos, vale aproximadamente:
\begin{multicols}{2}
\begin{choices}
	\correctchoice{$-1.9$ mm}
	\wrongchoice{$-3.1$ mm}
	\wrongchoice{$-1.0$ mm}
	\wrongchoice{$-0.4$ mm}
\end{choices}
\end{multicols}
\end{question}
}
%%%%%%%%%%%%%%%%%%%%%%%%%%%%%%%%%%%%%%%%%%%%%%%%%%%%%%%%%%%%%%%%%%%%%%%%%%%%%
\element{test1}{
\begin{question}{p6}
El valor del desplazamiento en dirección $X$ del nodo $A$,
calculado con los elementos de modos incompatibles, vale aproximadamente:
\begin{multicols}{2}
\begin{choices}
	\correctchoice{$-1.9$ mm}
	\wrongchoice{$-3.1$ mm}
	\wrongchoice{$-1.0$ mm}
	\wrongchoice{$-0.4$ mm}
\end{choices}
\end{multicols}
\end{question}
}
%%%%%%%%%%%%%%%%%%%%%%%%%%%%%%%%%%%%%%%%%%%%%%%%%%%%%%%%%%%%%%%%%%%%%%%%%%%%%
\element{test1}{
\begin{question}{p7}
Para evaluar la posible rotura del perfil, se visualizan 
los contornos de la tensión de Von Mises ({\tt pstr,6}).
Para los elementos isoparamétricos, el valor máximo es aproximadamente:
\begin{multicols}{2}
\begin{choices}
	\correctchoice{$47.7$ MPa}
	\wrongchoice{$26.1$ MPa}
	\wrongchoice{$58.2$ MPa}
	\wrongchoice{$34.6$ MPa}
\end{choices}
\end{multicols}
\end{question}
}
%%%%%%%%%%%%%%%%%%%%%%%%%%%%%%%%%%%%%%%%%%%%%%%%%%%%%%%%%%%%%%%%%%%%%%%%%%%%%
\element{test1}{
\begin{question}{p8}
Para evaluar la posible rotura del perfil, se visualizan 
los contornos de la tensión de Von Mises ({\tt pstr,6}).
Para los elementos con formulación mixta, el valor máximo es aproximadamente:
\begin{multicols}{2}
\begin{choices}
	\correctchoice{$50.2$ MPa}
	\wrongchoice{$29.4$ MPa}
	\wrongchoice{$63.1$ MPa}
	\wrongchoice{$38.2$ MPa}
\end{choices}
\end{multicols}
\end{question}
}
%%%%%%%%%%%%%%%%%%%%%%%%%%%%%%%%%%%%%%%%%%%%%%%%%%%%%%%%%%%%%%%%%%%%%%%%%%%%%
\element{test1}{
\begin{question}{p9}
Para evaluar la posible rotura del perfil, se visualizan 
los contornos de la tensión de Von Mises ({\tt pstr,6}).
Para los elementos con deformaciones mejoradas supuestas, el valor máximo es
aproximadamente:
\begin{multicols}{2}
\begin{choices}
	\correctchoice{$53.2$ MPa}
	\wrongchoice{$27.4$ MPa}
	\wrongchoice{$34.9$ MPa}
	\wrongchoice{$48.6$ MPa}
\end{choices}
\end{multicols}
\end{question}
}
%%%%%%%%%%%%%%%%%%%%%%%%%%%%%%%%%%%%%%%%%%%%%%%%%%%%%%%%%%%%%%%%%%%%%%%%%%%%%
\element{test1}{
\begin{question}{p10}
De acuerdo con los valores de los contornos de la tensión de
Von Mises, la posible zona de rotura del perfil es:
\begin{multicols}{2}
\begin{choices}
        \correctchoice{Con los tres modelos, en los puntos en que se aplican las cargas}
        \wrongchoice{Con el modelo de elementos con deformaciones mejoradas
supuestas, en los puntos en que se aplican las cargas. Con los modelos de
elementos mixtos e isoparamétricos, en la zona central de la cara interna}
        \wrongchoice{Con los modelos de elementos mixtos e isoparamétricos,
en los puntos en que se aplican las cargas. Con el modelo de elementos con
deformaciones mejoradas supuestas, en la zona central de la cara interna}
        \wrongchoice{Con los tres modelos, en la zona central de la cara
interna}
\end{choices}
\end{multicols}
\end{question}
}
%%%%%%%%%%%%%%%%%%%%%%%%%%%%%%%%%%%%%%%%%%%%%

\scoringDefaultS{b=1,m=-1/(N-1)}

\onecopy{1}{    

%%% beginning of the test sheet header:    

\noindent{\large\bf Método de los Elementos Finitos  \hfill MUECYM \hfill TEST \# 4}

\vspace*{.5cm}
\begin{minipage}{.4\linewidth}
  \centering 11 nov 2016
\end{minipage}

\begin{center}\em
Tiempo: 60 minutos.

  %Está prohibido el uso de teléfonos móviles.

  %Se atribuirá puntuación negativa a las respuestas incorrectas.

\end{center}
\vspace{1ex}

%%% end of the header

\shufflegroup{test1}
\insertgroup{test1}

\AMCcleardoublepage    
%\clearpage

\AMCformBegin    

%%% beginning of the answer sheet header

\noindent\AMCcode{nummat}{3}\hspace*{\fill}
\begin{minipage}{.7\linewidth}
$\longleftarrow{}$ Escriba su número de matrícula marcando los dígitos
en los recuadros (con ceros a la izquierda si el número es de menos de tres dígitos) y el nombre y apellidos debajo.

\vspace{3ex}

\namefield{\fbox{
   \begin{minipage}{.9\linewidth}
     Apellidos, Nombre:

     \vspace*{.5cm}\dotfill
     \vspace*{1mm}
   \end{minipage}
 }}
\end{minipage}

\begin{center}
 \bf\em Debe dar las respuestas exclusivamente en esta hoja (las respuestas en las demás hojas no serán tenidas en cuenta).
\end{center}

%%% end of the answer sheet header


\AMCform    

\AMCcleardoublepage    

}  

\end{document}
